\section{October 27}

We will talk about suprema of random processes. We proved a bunch of random results saying where a function of random variables are close to its expectation $f(X_1, \ldots, X_n) \approx \mathbb{E}[f]$, but actully how large is the expectation? In general, bounding the expectation $\mathbb{E}[f]$ requires case-specific arguments. It can't be easy since then, we would know all there is about the RV, and it would violate Murphy's law(?). 

We will develop general bounds for 
\begin{equation}
	f(Xt) = \sup_{t \in T} X_t
\end{equation}
Why is this important? 
\begin{enumerate}
	\item \textit{Norms}. For $x \in \mathbb{R}^n$, 
		\begin{equation}
			\|X\|_p = \sup_{\|v\|_q = 1} \langle u, X \rangle, \quad \frac{1}{p} + \frac{1}{q} = 1
		\end{equation}
		Similarly, if we have a matrix
		\begin{equation}
			\|X\|_{\mathrm{op}} = \sup_{\|u\|, \|v\| = 1} \langle X u, v \rangle
		\end{equation}

	\item \textit{Random Optimization}. Slightly more generally, let $T \subset \mathbb{R}^n$, $X \in \mathbb{R}^n$, and we would like to find 
		\begin{equation}
			\sup_{t \in T} \langle t, X \rangle
		\end{equation}
		This has applications in max cut, TSP, etc. 

	\item \textit{Empirical Processes}. The index set is a function class $\mathcal{F} = \{f : \mathcal{X} \to \mathbb{R} \}$. Assume we have some RV $X \in \mathcal{X}$ with $X_1, \ldots, X_n \sim X$. We care about 
		\begin{equation}
			\sup_{f \in \mathcal{F}} \frac{1}{n} \sum_i f(X_i) - \mathbb{E}[f(X)]
		\end{equation}
\end{enumerate}

The general principle is the following. If random process $t \to X_t$ is ``sufficiently continuous,'' then $\sup_t X_t$ is controlled by the ``geometric complexity'' of the index set $T$. Let's start with the simplest case. We have shown that if $X_1, \ldots, X_n \sim N(0, 1)$ iid, then $\sup{X_i} \sim \sqrt{\log{n}}$. We will generalize this to arbitrary subgaussian RVs. 

\begin{theorem}
	Suppose $|T| < +\infty$, and such $X_t$ is $\sigma^2$-subgaussian with $\mathbb{E}[X_t] = 0$. Then, 
	\begin{equation}
		\mathbb{E}[\sup_{t \in T} X_t] \leq \sqrt{2 \sigma^2 \log{|T|}}
	\end{equation}
	This upper bound is a very simple description of the complexity of $T$---it is the size. 
\end{theorem}
\begin{proof}
	For $\theta > 0$, we can write $X_t = \frac{1}{\theta} \log{e^{\theta X_t}}$. Then, we can bound by subgaussianness
	\begin{align}
		\sup_{t} X_t & = \frac{1}{\theta} \log{\ \sup_{t} e^{\theta X_t}} \\ 
								 & \leq \frac{1}{\theta} \log \sum_{t \in T} e^{\theta X_t}
	\end{align} 
	By Jensen's inequality, 
	\begin{align}
		\mathbb{E}[\sup_{t \in T} X_t] & \leq \frac{1}{\theta} \log \mathbb{E} \bigg[ \sum_{t \in T} e^{\theta X_t} \bigg] \\ 
																	 & \leq \frac{1}{\theta} \log \big( |T| e^{\theta^2 \sigma^2 /2} \big) \\ 
																	 & = \frac{1}{\theta} \big( \log{|T|} + \frac{\theta^2 \sigma^2}{2} \big)
	\end{align}
	So, the idea is that we want to transform this random variable to bring out the subgaussian-anity. 
\end{proof}

Is this a sharp bound? No, since the summation wastes a lot of space. Indeed, if $X_t$ were all dependent, i.e. $X_t$ were all the same, then 
\begin{equation}
	\mathbb{E}[\sup{t \in T} X_t] = \mathbb{E}[X] = 0
\end{equation}
However, if $X_t$ are all independent $N(0, \sigma^2)$, then this is correct. An exercise is that \begin{equation}
	\frac{\sup_{t \in T} X_t}{\sqrt{2 \sigma^2 \log{n}}} \to 1
\end{equation}
in probability. 

For infinite index sets, we must impose some dependence among the $X_t$. Another exercise: If $X_t$ are iid and unbounded with $T$ infinite, then $\sup_{t} X_t = +\infty$ a.s. To put more structure, we can posulate that if $t \approx s$, then $X_t \approx X_s$. To be precise, we will endow the index set $T$ with a metric and assume that $T$ comes with a metric. To quantify what $X_t \approx X_s$, we start with the simple definition. 

\begin{definition}[Lipschitz Process]
  $(X_t)$ is \textbf{Lipschitz} if there exists a (potentially random) Lipschitz constant $L > 0$ s.t. 
	\begin{equation}
	  |X_t - X_s| \leq L d(t, s)
	\end{equation}
	This is meaningful if $\mathbb{E}[L] < +\infty$. 
\end{definition}

\begin{example}
  Let $X_t = \langle t, g \rangle$ where $g \sim N(0, I)$. Then 
	\begin{equation}
	  |X_t - X_s| = |\langle t - s, g \rangle| \leq \|g\|_2 \cdot \|t - s\|_2
	\end{equation}
\end{example}

So how do we control the supremum of a Lipschitz process? The idea is the following. We first pick some finitely many points $t_1, \ldots, t_N \in T$ s.t. every other point of $t$ is close to $t_i$. Therefore, 
\begin{equation}
	\sup_{t \in T} X_t \lesssim \sup_{i = 1, \ldots, N} X_{t_i} \lesssim \sqrt{\log{n}}
\end{equation}

\begin{definition}[Covering Number]
	Given metric space $(T, d)$ and radius $\epsilon > 0$, a subset $N \subset T$ is called an $\epsilon$-net if for every $t \in T$, there exists a $\pi(t) \in N$ s.t. $d(x, \pi(t)) < \epsilon$. Equivalently, $d(t, N) = \inf_{s \in N} d(t, s) < \epsilon$. 

	The $\epsilon$-covering number of $(T, d)$ is 
	\begin{equation}
		N(T, d, \epsilon) \coloneqq \inf \big\{ |N| \colon N \text{ is an } \epsilon \text{-net} \big\}
	\end{equation}
\end{definition}

It turns out that these covering numbers are the geometric complexity that we will need. 

\begin{theorem}
  If $(X_t)$ is an $L$-lipschitz process, and if each $X_t$ is $\sigma^2$-subgaussian, then 
	\begin{equation}
		\mathbb{E}[ \sup_{t \in T} X_t] \leq \inf_{\epsilon > 0} \big\{ \epsilon \mathbb{E}[L] + \sqrt{2 \sigma^2 \log{N}} \big\}
	\end{equation}
	The $\epsilon \mathbb{E}[L]$ has to do with how much you lose by having to go to a representative point, and we see the familiar square root log. 
\end{theorem}
\begin{proof}
  Fix $\epsilon > 0$. Let $N$ be an $\epsilon$-net. For each $t \in T$, choose $\pi(t) \in N$ s.t. $d(\pi(t), t) < \epsilon$. Then 
	\begin{align}
		\sup_{t} X_t & = \sup_{t} \big\{ X_t - X_{\pi(t)} + X_{\pi(t)} \big\} \\ 
								 & \leq\sup_t (X_t - X_{\pi(t)}) + \sup_{t \in T} X_{\pi(t)} \\ 
								 & \leq L \sup_{t} d(t, \pi(t)) + \sup_{t \in N} X_t \\ 
								 & \leq L \epsilon + \sqrt{2 \sigma^2 \log{|N|}}
	\end{align}
	Now we take expectation and minimize over the $\epsilon > 0$-net. 
\end{proof}

A few remarks. First, this is interesting only if covering numbers are finite. This is not an artifact of the proof, but we really need finiteness. Second, $\inf$ over $\epsilon$ is a tradeoff between explaining ``independence'' and ``continuity'' at scale $\epsilon > 0$. (An exercise, $\epsilon \mapsto N(T, d, \epsilon)$ is decreasing.) Third, we didn't actually use our assumption on independence. 

So all of this comes down to understanding covering numbers. So how do we upper bound these covering numbers? The intuitive idea is that if we want the smallest cover, then we want the balls to be spread out as much as possible over representative points. This is called the \textit{packing number}. 

\begin{definition}[Packing Number]
  If $(T, d)$ is a metric space and $\epsilon > 0$, then a subset $P \subset T$ is an \textbf{$\epsilon$-packing} if for all $t, s \in P$ with $t \neq s$, $d(t, s) > \epsilon$. The $\epsilon$-packing number is defined 
	\begin{equation}
		P(T, d, \epsilon) = \max \{ |P| \colon P \text{ is an } \epsilon \text{-packing}\}
	\end{equation}
	Note that this may not exist for some spaces, and often we will assume compactness. 
\end{definition}

Another exercise is that $\epsilon \to fP(T, d, \epsilon)$ is decreasing. 

\begin{theorem}[Covering-Packing Duality]
  For all $\epsilon > 0$. 
	\begin{equation}
	  P(T, d, 2 \epsilon) \leq N(T, d, \epsilon) \leq P(T, d, \epsilon)
	\end{equation}
\end{theorem}
\begin{proof}
  Say $ \subset T$ is a $2\epsilon$-packing and $N \subset T$ is an $\epsilon$-net. For any $t \in T$, let $\pi(t) \in N$ s.t. $d(t, \pi(t)) \leq \epsilon$. We claim that $\pi: P \to N$ is injective. For any $t \neq s$ in $P$, 
	\begin{equation}
		2 \epsilon < d(t, s) \leq d(t, \pi(t)) + d(\pi(t), \pi(s)) + d(\pi(s), s) \leq 2 \epsilon + d(\pi(t), \pi(s))
	\end{equation}
	So, $d(\pi(t), \pi(s)) > 0$. So, $|P| \leq |N|$. 

	For the other inequality, let $P \subset T$ be a maximal $\epsilon$-packing. We claim that $P$ is an $\epsilon$-net. If not, there exists $t \in T \setminus P$ s.t. $d(t, s) > \epsilon$ for all $s \in P$. Then, $P^1 = P \cup \{t\}$ is an $\epsilon$-packing, which means $|P^1| = |P| + 1$, contradiction. 
\end{proof}

\begin{theorem}[Volumetric Bound]
	For $A \subset \mathbb{R}^n$ with $2$-norm, 
	\begin{equation}
		\frac{\mathrm{vol}(A)}{\mathrm{vol}(\epsilon B_n)} \leq N(A, \|\cdot \|_2, \epsilon) \leq P(A, \|\cdot\|_2, \epsilon) \leq \frac{\mathrm{vol}(A + \frac{\epsilon}{2} B^n_1)}{\mathrm{vol} \frac{\epsilon}{2} B_1^n}
	\end{equation}
	where $B_1^n$ is the ball of radius $1$ in $\mathbb{R}^n$, and 
	\begin{equation}
		A + \frac{\epsilon}{2} B_n^1 = A_{\epsilon/2} = \{ y \in \mathbb{R}^n \colon \exists x \in A, \|x - y\| \leq \epsilon/2 \}
	\end{equation}
\end{theorem}
\begin{proof}
	Let $N \subset A$ be an $\epsilon$-net. Then, $A \subset \cup_{x \in N} (x + \epsilon B_1^n)$, which impolies 
	\begin{equation}
		\mathrm{vol}(A) \leq |N| \mathrm{vol} \big( \sum B_1^n \big) \implies |N| \geq \frac{\mathrm{vol}(A)}{\mathrm{vol} \big( \sum B_1^n \big)}
	\end{equation}
	If $f \subset A$ is an $\epsilon$ packing, then
	\begin{equation}
		\bigsqcup_{x \in P} \big( x + \frac{\epsilon}{2} B_1^n \big) \subset A_{\epsilon/2} \implies \underbrace{\sum_{x \in P} \mathrm{vol} \big( x + \frac{\epsilon}{2} B_1^n \big)}_{|P| \mathrm{vol}( \frac{\epsilon}{2} B_1^n)} \leq \mathrm{vol}(A_{\epsilon/2})
	\end{equation}
\end{proof}

As a corollary, we have precise bounds on the covering number of certain sets, like the Euclidean ball. 

\begin{corollary}
  If $A = B_1^n$, then 
	\begin{equation}
		\bigg(\frac{1}{\epsilon} \bigg)^n \leq N(A, \|\cdot\|_2, \epsilon) \leq \bigg( \frac{3}{\epsilon} \bigg)^n
	\end{equation}
\end{corollary}

Remark. If $A$ is bounded and has nonempty interior (and nice in the sense as in no cantor sets, etc.), then $N(A, \|\cdot\|, \epsilon) \sim \epsilon^{-n}$. This can actually be used to define dimension, such as the \textit{metric dimension} of a metric space. 

\begin{example}
  Listed. 
	\begin{enumerate}
	  \item If $g = (g_1, \ldots, g_n)$ are independent and $\sigma^2$-subgaussian, 
			\begin{equation}
				T = B_1^n, \quad X_t = \langle g, t \rangle,
			\end{equation}
			We can upper bound 
			\begin{equation}
				\sup_{t \in T} X_t = \sup_{\|t\|_2 \leq 1} \langle g, t \rangle = \|g\|_2 
			\end{equation}

			As an exercise,show that $X_t$ is $\sigma^2$-subgaussian for all $t \in R_1^n$. Moreover, $|X_t - X_s| \leq \|g_2\| \|t - s\|_2$, so it is Lipschitz. For all $\epsilon > 0$, 
			\begin{align}
				\mathbb{E}[\|g\|_2] & \leq \epsilon \mathbb{E}[\|g\|_2] + \sqrt{2 \sigma^2 \log{N(B_1^n , \|\cdot\|_2, \epsilon)}} \\ 
														& \leq \epsilon \mathbb{E}[\|g\|_2] + \sqrt{2 \sigma^2 n \log(3/\epsilon)} \\ 
				\mathbb{E}[\|g\|_2] & \leq \inf_{\epsilon > 0} \sqrt{2 \sigma^2 n} \frac{\sqrt{\log{3/\epsilon}}}{1 - \epsilon} 
			\end{align}
			and therefore $\mathbb{E}[\|g\|_2] \leq c \sqrt{\sigma^2 n} \approx 2.035$. 

		\item \textit{Maximum Eigenvalues}> If $G$ is a matrix with entries $G_{ij} = G_{ji}$ independent $\sigma^2$-subgaussians. Let 
			\begin{equation}
				X_t = |v^T G v|, \quad T = B_1^n, \quad \sup_{t \in T} X_t = \sup_{\|v\|_2 = 1} | v G v| = \|G\|_{\mathrm{op}}
			\end{equation}
			As an exercise, show that $X_v$ is $2\sigma^2$-subgaussian. We see that 
			\begin{equation}
				|X_v - X_u| = |v^T G v - u^T G u| = | v^T G(V - u) + (u - v)^T G u| \leq 2 \|G\|_{\mathrm{op}} |v - u|_2
			\end{equation}
			and therefore, 
			\begin{equation}
				\mathbb{E}{\|G\|_{\mathrm{op}}} \leq \inf_{\epsilon > 0} \frac{\sqrt{4 \sigma^2 n \log (3/\epsilon)}}{1 - 2 \epsilon} \approx 2.25 \cdot 2 \sqrt{\sigma^2 n} \sim O(\sqrt{n})
			\end{equation}
	\end{enumerate}
\end{example}
